% Generated from project outputs. Do not edit by hand.
\begin{table}[!htbp]
\centering
\begin{threeparttable}
\caption{Capacity-curve sensitivity: age 2016-2019 fixed\_baseline}
\label{tab:robustness-age-2016-2019-fixed-baseline}
\small
\setlength{\tabcolsep}{4pt}
\renewcommand{\arraystretch}{1.15}
\begin{tabularx}{\linewidth}{@{}Xrrrrr@{}}
\toprule
Variant & Span & Q500 & Q1500 & Q3000 & Q5000 \\
\midrule
Main (14-60) & 0.300 & 667.5 & 1,774.5 & 3,148.0 & 5,403.0 \\
v1 (13-59) & 0.300 & 670.5 & 1,781.5 & 3,145.1 & 5,384.6 \\
v2 (15-61) & 0.300 & 688.2 & 1,778.8 & 3,156.9 & 5,402.7 \\
v3 (12-58) & 0.300 & 669.8 & 1,781.3 & 3,145.1 & 5,388.3 \\
v4 (16-62) & 0.300 & 774.0 & 1,778.4 & 3,156.9 & 5,426.8 \\
\bottomrule
\end{tabularx}
\begin{tablenotes}[flushleft]
\footnotesize
\item Entries are predicted income at t+1 (quetzales); column headings give income at t. General sample, labor income, population-weighted LOESS. -- indicates no supported prediction. Retuned selects a span for each cutoff; fixed\_baseline holds the main-sample span fixed.
\end{tablenotes}
\end{threeparttable}
\end{table}
